% A1
\subsection{Step-Level Breakdown Evidence}\label{app:A1}

This appendix supports step-level claims about where AWM gains and
losses come from in \texttt{C1}. The evidence below is drawn from the
audited step-level decomposition of the \texttt{C1} sites.

\subsubsection{Representative step-level deltas from the source
output}

{\def\LTcaptype{} % do not increment counter
\begin{longtable}[]{@{}>{\raggedright\arraybackslash}p{0.14\linewidth}>{\raggedright\arraybackslash}p{0.22\linewidth}>{\raggedleft\arraybackslash}p{0.14\linewidth}>{\raggedleft\arraybackslash}p{0.14\linewidth}>{\raggedleft\arraybackslash}p{0.14\linewidth}>{\raggedleft\arraybackslash}p{0.18\linewidth}@{}}
\toprule\noalign{}
Website & Condition & \texttt{CLICK\ dElem\ /\ dStepSR} &
\texttt{TYPE\ dElem\ /\ dActF1\ /\ dStepSR} &
\texttt{first\_half\ dStepSR} & \texttt{second\_half\ dStepSR} \\
\midrule\noalign{}
\endhead
\bottomrule\noalign{}
\endlastfoot
\texttt{budget} & \texttt{offline\_wf\ -\ no\_workflow} &
\texttt{-12.1\ /\ -12.1} & \texttt{+0.0\ /\ +8.4\ /\ +9.1} &
\texttt{-1.9} & \texttt{-10.9} \\
\texttt{kayak} & \texttt{offline\_wf\ -\ no\_workflow} &
\texttt{+10.7\ /\ +10.7} & \texttt{+0.0\ /\ +1.4\ /\ +0.0} &
\texttt{+0.0} & \texttt{+13.0} \\
\texttt{newegg} & \texttt{offline\_wf\ -\ no\_workflow} &
\texttt{+9.1\ /\ +6.1} & \texttt{+33.3\ /\ +50.0\ /\ +33.3} &
\texttt{+4.4} & \texttt{+7.1} \\
\texttt{sixflags} & \texttt{offline\_wf\ -\ no\_workflow} &
\texttt{-6.1\ /\ -6.1} & \texttt{+0.0\ /\ +0.0\ /\ +0.0} & \texttt{-5.9}
& \texttt{-3.3} \\
\texttt{united} & \texttt{offline\_wf\ -\ no\_workflow} &
\texttt{+3.6\ /\ +3.6} & \texttt{+16.7\ /\ +2.8\ /\ +0.0} &
\texttt{+3.0} & \texttt{+0.0} \\
\texttt{yellowpages} & \texttt{offline\_wf\ -\ no\_workflow} &
\texttt{+0.0\ /\ +0.0} & \texttt{+28.6\ /\ +19.5\ /\ +28.6} &
\texttt{+0.0} & \texttt{+7.1} \\
\end{longtable}
}

\subsubsection{Source excerpts}

\paragraph{Kayak}

\begin{verbatim}
[Delta: offline_wf - no_workflow] by action_type
  CLICK                   +10.7%     +3.6%    +10.7%
  SKIP                     +0.0%     +0.0%     +0.0%
  TYPE                     +0.0%     +1.4%     +0.0%

[Delta: offline_wf - no_workflow] by step_half
  first_half               +0.0%     +0.3%     +0.0%
  second_half             +13.0%     +4.3%    +13.0%
\end{verbatim}

\paragraph{Budget}

\begin{verbatim}
[Delta: offline_wf - no_workflow] by action_type
  CLICK                   -12.1%     -1.7%    -12.1%
  SELECT                   +0.0%     +0.0%     +0.0%
  SKIP                     +0.0%     +0.0%     +0.0%
  TYPE                     +0.0%     +8.4%     +9.1%

[Delta: offline_wf - no_workflow] by step_half
  first_half               -3.8%     +3.6%     -1.9%
  second_half             -10.9%     -4.3%    -10.9%
\end{verbatim}

\paragraph{Newegg}

\begin{verbatim}
[Delta: offline_wf - no_workflow] by action_type
  CLICK                    +9.1%     -3.0%     +6.1%
  SELECT                   +0.0%    +20.0%    +25.0%
  SKIP                     +0.0%     +0.0%     +0.0%
  TYPE                    +33.3%    +50.0%    +33.3%
\end{verbatim}

\subsubsection{Interpretation limits}

\begin{itemize}
\tightlist
\item
  This appendix supports direction-of-effect claims, not causal proof.
\item
  \texttt{TYPE} counts are small on several sites, so the main robust
  use is qualitative: \texttt{TYPE} gains are often non-negative while
  \texttt{CLICK} varies sharply by site.
\item
  Later-step accumulation is strongest on \texttt{kayak}; other sites
  show weaker or mixed support.
\end{itemize}


% A2
\clearpage
\subsection{Paired-Case Summary Evidence}\label{app:A2}

This appendix supports claims about how often workflow actually changes
behavior, and whether changed steps are net positive or net negative.
The evidence below summarizes the audited paired-case distributions over
the seven \texttt{C1} sites.

\subsubsection{\texorpdfstring{A2.1 Corrected overall totals for the seven
\texttt{C1}
websites}{A2.1 Corrected overall totals for the seven C1 websites}}

This table corrects the earlier erroneous denominator \texttt{523}.\\
The total paired-step count for the seven audited \texttt{C1} sites is
\texttt{475}.

{\def\LTcaptype{} % do not increment counter
\begin{longtable}[]{@{}>{\raggedright\arraybackslash}p{0.28\linewidth}>{\raggedright\arraybackslash}p{0.46\linewidth}>{\raggedright\arraybackslash}p{0.22\linewidth}@{}}
\toprule\noalign{}
Category & Count & Pct of all paired steps \\
\midrule\noalign{}
\endhead
\bottomrule\noalign{}
\endlastfoot
\texttt{positive} & 26 & \texttt{5.5\%} \\
\texttt{negative} & 24 & \texttt{5.1\%} \\
\texttt{ineffective} & 252 & \texttt{53.1\%} \\
\texttt{redundant} & 173 & \texttt{36.4\%} \\
\texttt{changed\ =\ positive\ +\ negative} & 50 & \texttt{10.5\%} \\
\end{longtable}
}

Among changed steps only:

{\def\LTcaptype{} % do not increment counter
\begin{longtable}[]{@{}>{\raggedright\arraybackslash}p{0.36\linewidth}>{\centering\arraybackslash}p{0.60\linewidth}@{}}
\toprule\noalign{}
Metric & Value \\
\midrule\noalign{}
\endhead
\bottomrule\noalign{}
\endlastfoot
\texttt{positive\ /\ (positive\ +\ negative)} &
\texttt{26\ /\ 50\ =\ 52.0\%} \\
\texttt{net\ gain} & \texttt{+2\ steps} \\
\end{longtable}
}

\subsubsection{\texorpdfstring{A2.2 Site-wise evidence from
\texttt{C1}}{A2.2 Site-wise evidence from C1}}

{\def\LTcaptype{} % do not increment counter
\begin{longtable}[]{@{}>{\raggedright\arraybackslash}p{0.12\linewidth}>{\raggedleft\arraybackslash}p{0.20\linewidth}>{\raggedleft\arraybackslash}p{0.11\linewidth}>{\raggedleft\arraybackslash}p{0.11\linewidth}>{\raggedleft\arraybackslash}p{0.11\linewidth}>{\raggedleft\arraybackslash}p{0.11\linewidth}>{\raggedright\arraybackslash}p{0.20\linewidth}@{}}
\toprule\noalign{}
Website & Total & Positive & Negative & Ineffective & Redundant &
Verdict pattern \\
\midrule\noalign{}
\endhead
\bottomrule\noalign{}
\endlastfoot
\texttt{budget} & 99 & 6 & 12 & 53 & 28 & negative-dominated \\
\texttt{kayak} & 48 & 3 & 0 & 22 & 23 & reproduced, no negatives \\
\texttt{kohls} & 53 & 2 & 2 & 33 & 16 & mixed \\
\texttt{newegg} & 87 & 5 & 0 & 56 & 26 & reproduced, no negatives \\
\texttt{sixflags} & 64 & 3 & 6 & 23 & 32 & negative-dominated \\
\texttt{united} & 63 & 3 & 2 & 35 & 23 & slight positive edge \\
\texttt{yellowpages} & 61 & 4 & 2 & 30 & 25 & slight positive edge \\
\end{longtable}
}

\subsubsection{Source excerpts}

\paragraph{Kayak}

\begin{verbatim}
[Overall Distribution]
  positive             3    6.2%
  negative             0    0.0%
  ineffective         22   45.8%
  redundant           23   47.9%

[Workflow Impact]
  Positive / (Positive + Negative) = 3/3 = 100.0%
  Net gain (positive - negative)   = +3 steps
\end{verbatim}

\paragraph{Newegg}

\begin{verbatim}
[Overall Distribution]
  positive             5    5.7%
  negative             0    0.0%
  ineffective         56   64.4%
  redundant           26   29.9%
\end{verbatim}

\paragraph{Budget}

\begin{verbatim}
[Overall Distribution]
  positive             6    6.1%
  negative            12   12.1%
  ineffective         53   53.5%
  redundant           28   28.3%
\end{verbatim}

\paragraph{Sixflags}

\begin{verbatim}
[Overall Distribution]
  positive             3    4.7%
  negative             6    9.4%
  ineffective         23   35.9%
  redundant           32   50.0%
\end{verbatim}

\subsubsection{Interpretation limits}

\begin{itemize}
\tightlist
\item
  This appendix supports statements about behavioral footprint and sign
  of net change.
\item
  It should not be used to claim that all changed steps are equally
  important.
\item
  The corrected denominator \texttt{475} should be used consistently
  wherever these \texttt{C1} distributions are discussed.
\end{itemize}


% A3
\clearpage
\subsection{First-Run Target-Site Result Note}\label{app:A3}

This appendix records the first-run result facts on the two target
sites and preserves the corresponding diagnostic output as exploratory
tracing material. The note below is based on the audited target-site
error summaries and their associated negative-case review.

\subsubsection{Summary table}

{\def\LTcaptype{} % do not increment counter
\begin{longtable}[]{@{}>{\raggedright\arraybackslash}p{0.12\linewidth}>{\raggedright\arraybackslash}p{0.20\linewidth}>{\raggedleft\arraybackslash}p{0.11\linewidth}>{\raggedleft\arraybackslash}p{0.11\linewidth}>{\raggedleft\arraybackslash}p{0.11\linewidth}>{\raggedleft\arraybackslash}p{0.11\linewidth}>{\raggedright\arraybackslash}p{0.20\linewidth}@{}}
\toprule\noalign{}
Target & Model / setup & Positive & Negative & Skip-rate shift vs source
& CLICK share of negatives & Safe reading \\
\midrule\noalign{}
\endhead
\bottomrule\noalign{}
\endlastfoot
\texttt{tripadvisor} & \texttt{qwen/qwen3.5-397b-\allowbreak a17b\ /\ online\_wf} &
4 & 18 & \texttt{+16pp} (\texttt{44.9\%} vs \texttt{29.2\%}) &
\texttt{78\%} & first-run underperformance with weaker baseline
candidate quality \\
\texttt{reddit} & \texttt{qwen/qwen3.5-397b-\allowbreak a17b\ /\ online\_wf} & 2 & 5
& \texttt{+4.7pp} (\texttt{33.9\%} vs \texttt{29.2\%}) & \texttt{80\%} &
first-run underperformance without obvious baseline collapse \\
\end{longtable}
}

\subsubsection{Source excerpts}

\paragraph{Reddit}

\begin{verbatim}
[Degradation Profile: online_wf]
  Positive: 2, Negative: 5, Ineffective: 87, Redundant: 86

[Hypothesis B: Element grounding collapse]
  Skip rate (target baseline): 33.9%
  Skip rate (source baseline): 29.2%
  CLICK Elem Acc (target baseline): 84.2%
  CLICK Elem Acc (source baseline): 75.0%
  -> No strong B evidence (baseline performs similarly)

[Hypothesis A: Workflow content mismatch]
  Negative steps: 5
    CLICK: 4  TYPE: 1
    CLICK share of negatives: 80%
  -> A evidence: Negatives dominated by CLICK (element-level mismatch, not value mismatch)

[Preliminary Verdict]
  -> Primarily Hypothesis A: workflow content is mismatched
\end{verbatim}

\paragraph{Tripadvisor}

\begin{verbatim}
[Degradation Profile: online_wf]
  Positive: 4, Negative: 18, Ineffective: 131, Redundant: 72

[Hypothesis B: Element grounding collapse]
  Skip rate (target baseline): 44.9%
  Skip rate (source baseline): 29.2%
  CLICK Elem Acc (target baseline): 81.2%
  CLICK Elem Acc (source baseline): 75.0%
  -> B evidence: Skip rate +16pp higher on target site

[Hypothesis A: Workflow content mismatch]
  Negative steps: 18
    CLICK: 14  TYPE: 4
    CLICK share of negatives: 78%
  -> A evidence: Negatives dominated by CLICK (element-level mismatch, not value mismatch)

[Preliminary Verdict]
  -> Mixed: both element grounding degradation AND workflow mismatch
\end{verbatim}

\subsubsection{Interpretation limits}

\begin{itemize}
\tightlist
\item
  This appendix is best used as first-run result tracing, not as a
  report-level mechanism proof.
\item
  The diagnostic labels in the underlying output are useful for
  internal inspection, but they should not be promoted to workflow
  provenance claims or full causal decomposition.
\item
  The strongest safe summary is:

  \begin{itemize}
  \tightlist
  \item
    \texttt{tripadvisor}: first-run underperformance plus weaker
    baseline candidate quality
  \item
    \texttt{reddit}: first-run underperformance without obvious
    baseline collapse
  \end{itemize}
\end{itemize}


% A4
\clearpage
\subsection{C4 Result Table}\label{app:A4}

This appendix supports \texttt{C4} claims about \texttt{NL\ vs\ HTML}
and \texttt{code\ vs\ text\ workflow}. The tables below consolidate the
audited first-run \texttt{C4} representation results.

\subsubsection{\texorpdfstring{A4.1 \texttt{NL\ vs\ HTML} first-run
results}{A4.1 NL vs HTML first-run results}}

{\def\LTcaptype{} % do not increment counter
\begin{longtable}[]{@{}>{\raggedright\arraybackslash}p{0.14\linewidth}>{\raggedright\arraybackslash}p{0.22\linewidth}>{\raggedleft\arraybackslash}p{0.14\linewidth}>{\raggedleft\arraybackslash}p{0.14\linewidth}>{\raggedleft\arraybackslash}p{0.14\linewidth}>{\raggedleft\arraybackslash}p{0.18\linewidth}@{}}
\toprule\noalign{}
Website & Condition & Element Acc & Action F1 & Step SR & SR \\
\midrule\noalign{}
\endhead
\bottomrule\noalign{}
\endlastfoot
\texttt{kayak} & \texttt{desc\_only} & 54.8 & 61.2 & 45.3 & 0.0 \\
\texttt{kayak} & \texttt{html\_only} & 52.4 & 63.5 & 48.0 & 0.0 \\
\texttt{kayak} & \texttt{desc\_html} & 54.8 & 63.5 & 50.3 & 0.0 \\
\texttt{united} & \texttt{desc\_only} & 57.9 & 61.9 & 55.9 & 33.3 \\
\texttt{united} & \texttt{html\_only} & 61.2 & 64.5 & 60.6 & 33.3 \\
\texttt{united} & \texttt{desc\_html} & 57.9 & 61.1 & 51.7 & 16.7 \\
\texttt{newegg} & \texttt{desc\_only} & 42.0 & 44.3 & 30.5 & 0.0 \\
\texttt{newegg} & \texttt{html\_only} & 36.2 & 42.6 & 23.9 & 0.0 \\
\texttt{newegg} & \texttt{desc\_html} & 38.0 & 42.4 & 30.3 & 0.0 \\
\end{longtable}
}

\subsubsection{\texorpdfstring{A4.2 \texttt{code\ vs\ text\ workflow}
first-run
result}{A4.2 code vs text workflow first-run result}}

{\def\LTcaptype{} % do not increment counter
\begin{longtable}[]{@{}>{\raggedright\arraybackslash}p{0.14\linewidth}>{\raggedright\arraybackslash}p{0.22\linewidth}>{\raggedleft\arraybackslash}p{0.14\linewidth}>{\raggedleft\arraybackslash}p{0.14\linewidth}>{\raggedleft\arraybackslash}p{0.14\linewidth}>{\raggedleft\arraybackslash}p{0.18\linewidth}@{}}
\toprule\noalign{}
Website & Condition & Element Acc & Action F1 & Step SR & SR \\
\midrule\noalign{}
\endhead
\bottomrule\noalign{}
\endlastfoot
\texttt{kayak} & \texttt{text\_wf} & 55.3 & 60.8 & 45.8 & 0.0 \\
\texttt{kayak} & \texttt{code\_wf} & 52.4 & 63.4 & 48.0 & 0.0 \\
\end{longtable}
}

\subsubsection{Interpretive summary}

The three-site \texttt{NL\ vs\ HTML} comparison does not support the
paper's directional claim on Mind2Web. Under the current prompt-level
code representation, the first \texttt{kayak} run suggests only modest
differences between \texttt{code\ workflow} and
\texttt{text\ workflow}.

\subsubsection{Interpretation limits}

\begin{itemize}
\tightlist
\item
  This appendix supports first-run result claims only.
\item
  \texttt{code\ vs\ text\ workflow} currently has one audited site
  (\texttt{kayak}), not a three-site result.
\end{itemize}


% A5
\clearpage
\subsection{C5 Workflow Quality Table}\label{app:A5}

This appendix supports \texttt{C5} claims about workflow compactness,
overlap, and prompt-level usage proxies. The table below summarizes the
audited first-run workflow-quality statistics under the current
approximation.

\subsubsection{First-run quality
table}

{\def\LTcaptype{} % do not increment counter
\begin{longtable}[]{@{}>{\raggedright\arraybackslash}p{0.12\linewidth}>{\raggedright\arraybackslash}p{0.20\linewidth}>{\raggedleft\arraybackslash}p{0.11\linewidth}>{\raggedleft\arraybackslash}p{0.11\linewidth}>{\raggedleft\arraybackslash}p{0.11\linewidth}>{\raggedleft\arraybackslash}p{0.11\linewidth}>{\raggedright\arraybackslash}p{0.20\linewidth}@{}}
\toprule\noalign{}
Website & Condition & \#workflows & coverage & function overlap &
utility rate & Current verdict \\
\midrule\noalign{}
\endhead
\bottomrule\noalign{}
\endlastfoot
\texttt{kayak} & \texttt{lm\_wf} & 7 & 1.0 & 0.0 & 1.0 &
\texttt{reproduced\ (first\ run)} \\
\texttt{newegg} & \texttt{lm\_wf} & 5 & 1.0 & 0.0333 & 1.0 &
\texttt{reproduced\ (first\ run)} \\
\texttt{united} & \texttt{lm\_wf} & 5 & 1.0 & 0.0 & 1.0 &
\texttt{reproduced\ (first\ run)} \\
\end{longtable}
}

\subsubsection{Interpretive summary}

Across the three audited sites, the first-run evidence supports the
directional workflow-quality conclusion under the current approximation:
the workflow libraries are compact, pairwise overlap is low, and the
prompt-level usage proxies are high. The main caveat is that
\texttt{utility\ rate} remains a broad proxy rather than a strict
adherence measure.

\subsubsection{Interpretation limits}

\begin{itemize}
\tightlist
\item
  \texttt{coverage\ =\ 1.0} and \texttt{utility\ rate\ =\ 1.0} here are
  prompt-level proxies.
\item
  These values should not be reworded as ``the model always followed the
  workflow.''
\item
  The safest wording is: the workflow text was consistently present and
  detectable in the prompt context under the current approximation.
\end{itemize}


% A6
\clearpage
\subsection{Alignment-Rate Note}\label{app:A6}

This appendix provides a traceable source block for the exploratory
alignment-rate numbers used in \texttt{§10.2}. The heuristic below is
reconstructed from the alignment-rate analysis used during the audit.

\subsubsection{What this metric is}

The alignment-rate audit computes a workflow-target alignment heuristic
using keyword overlap between:

\begin{itemize}
\tightlist
\item
  workflow pattern descriptions
\item
  target-element context extracted from the observation HTML
\end{itemize}

It is explicitly \textbf{not} a semantic-match metric.

\subsubsection{Current site-wise
values}

{\def\LTcaptype{} % do not increment counter
\begin{longtable}[]{@{}>{\raggedright\arraybackslash}p{0.28\linewidth}>{\raggedleft\arraybackslash}p{0.46\linewidth}>{\raggedright\arraybackslash}p{0.22\linewidth}@{}}
\toprule\noalign{}
Website & Alignment rate & C1 status \\
\midrule\noalign{}
\endhead
\bottomrule\noalign{}
\endlastfoot
\texttt{budget} & 97.2\% & not reproduced \\
\texttt{sixflags} & 94.2\% & not reproduced \\
\texttt{newegg} & 90.9\% & reproduced \\
\texttt{kayak} & 70.6\% & reproduced \\
\end{longtable}
}

\subsubsection{Safe interpretation}

\begin{itemize}
\tightlist
\item
  The metric is suitable for supporting the claim that \textbf{surface
  alignment does not reliably predict performance}.
\item
  It is \textbf{not} strong enough to support semantic-fit claims on its
  own.
\item
  \texttt{newegg} is an explicit exception to any monotonic ``higher
  alignment = worse performance'' story.
\end{itemize}

\subsubsection{Source excerpt}

\begin{verbatim}
Heuristic: keyword overlap between workflow pattern descriptions
           and observation element context for each target action.
Limitations: label-lag, keyword ambiguity, placeholder fallback,
             no value-level matching for TYPE/SELECT.

SUMMARY TABLE
  kayak               17      6     34       24    70.6%
  newegg              14      7     44       40    90.9%
  budget              35     10     71       69    97.2%
  sixflags            18      7     52       49    94.2%
\end{verbatim}


% A7
\clearpage
\subsection{Offline-vs-Online Trade-Off}\label{app:A7}

\subsubsection{Overview}

This appendix is retained as an exploratory rereading of existing
Mind2Web outputs. It should not be used as a primary final-report
appendix for target-site mechanism claims.

The goal here is not to introduce a new experiment. It is to re-read
existing outputs from a mechanism perspective and ask:

\begin{itemize}
\tightlist
\item
  what kinds of step-level gains does \texttt{offline\_wf} produce on
  \texttt{kayak}?
\item
  what kinds of step-level gains does \texttt{online\_wf} produce on the
  same site?
\item
  what first-run result pattern appears on the farther targets
  \texttt{tripadvisor} and \texttt{reddit}?
\end{itemize}

This appendix is based on four source types: audited step-level
summaries, audited paired-case summaries, workflow texts from
\texttt{kayak}, and manually verified negative cases from the
\texttt{tripadvisor} and \texttt{reddit} target-site runs.

\subsubsection{\texorpdfstring{A7.1 Same-Site Comparison on
\texttt{kayak}}{A7.1 Same-Site Comparison on kayak}}

\paragraph{Aggregate outcome}

From the audited step-level summary:

{\def\LTcaptype{} % do not increment counter
\begin{longtable}[]{@{}>{\raggedright\arraybackslash}p{0.20\linewidth}>{\raggedleft\arraybackslash}p{0.30\linewidth}>{\raggedleft\arraybackslash}p{0.22\linewidth}>{\raggedleft\arraybackslash}p{0.24\linewidth}@{}}
\toprule\noalign{}
condition & Elem Acc & Act F1 & Step SR \\
\midrule\noalign{}
\endhead
\bottomrule\noalign{}
\endlastfoot
no\_workflow & 52.1\% & 60.9\% & 47.9\% \\
offline\_wf & 58.3\% & 63.2\% & 54.2\% \\
online\_wf & 56.2\% & 65.4\% & 54.2\% \\
\end{longtable}
}

Observation:

\begin{itemize}
\tightlist
\item
  both \texttt{offline\_wf} and \texttt{online\_wf} improve
  \texttt{Step\ SR} by \texttt{+6.3pp} over baseline on \texttt{kayak}
\item
  \texttt{offline\_wf} has slightly higher \texttt{Elem\ Acc}
\item
  \texttt{online\_wf} has slightly higher \texttt{Act\ F1}
\end{itemize}

This already suggests that the two conditions are not helping in exactly
the same way.

\paragraph{Action-type
decomposition}

From the audited step-level summary:

{\def\LTcaptype{} % do not increment counter
\begin{longtable}[]{@{}>{\raggedright\arraybackslash}p{0.20\linewidth}>{\raggedleft\arraybackslash}p{0.30\linewidth}>{\raggedleft\arraybackslash}p{0.22\linewidth}>{\raggedleft\arraybackslash}p{0.24\linewidth}@{}}
\toprule\noalign{}
delta vs no\_workflow & CLICK dStepSR & TYPE dActF1 & TYPE dStepSR \\
\midrule\noalign{}
\endhead
\bottomrule\noalign{}
\endlastfoot
offline\_wf & +10.7\% & +1.4\% & +0.0\% \\
online\_wf & +7.1\% & +19.2\% & +16.7\% \\
\end{longtable}
}

Reading:

\begin{itemize}
\tightlist
\item
  \texttt{offline\_wf} contributes more on \texttt{CLICK} grounding
\item
  \texttt{online\_wf} contributes much more on \texttt{TYPE}
  formatting/value guidance
\end{itemize}

This is the clearest quantitative sign of a trade-off rather than a
simple dominance relation.

\paragraph{\texorpdfstring{Direct paired comparison:
\texttt{online\_wf} vs
\texttt{offline\_wf}}{Direct paired comparison: online\_wf vs offline\_wf}}

From the direct offline-vs-online paired comparison:

\begin{center}
\begin{tabular}{@{}>{\raggedright\arraybackslash}p{0.28\linewidth}>{\raggedright\arraybackslash}p{0.46\linewidth}>{\raggedright\arraybackslash}p{0.22\linewidth}@{}}
\toprule
category & count & pct \\
\midrule
positive & 1 & 2.1\% \\
negative & 1 & 2.1\% \\
ineffective & 21 & 43.8\% \\
redundant & 25 & 52.1\% \\
\bottomrule
\end{tabular}
\end{center}

By action type:

{\def\LTcaptype{} % do not increment counter
\begin{longtable}[]{@{}>{\raggedright\arraybackslash}p{0.20\linewidth}>{\raggedleft\arraybackslash}p{0.30\linewidth}>{\raggedleft\arraybackslash}p{0.22\linewidth}>{\raggedleft\arraybackslash}p{0.24\linewidth}@{}}
\toprule\noalign{}
type & pos & neg & total \\
\midrule\noalign{}
\endhead
\bottomrule\noalign{}
\endlastfoot
CLICK & 0 & 1 & 28 \\
TYPE & 1 & 0 & 6 \\
\end{longtable}
}

The only positive case for \texttt{online\_wf} over \texttt{offline\_wf}
is a TYPE-format correction:

\begin{itemize}
\tightlist
\item
  task: \texttt{Find\ a\ cheapest\ SUV\ in\ Brooklyn\ for\ 1\ day}
\item
  target: \texttt{TYPE\ {[}35686{]}\ {[}PENN\ STATION{]}}
\item
  offline-side prediction: \texttt{Penn\ Station}
\item
  online-side prediction: \texttt{PENN\ STATION}
\end{itemize}

The only negative case for \texttt{online\_wf} against
\texttt{offline\_wf} is a CLICK choice on the same task:

\begin{itemize}
\tightlist
\item
  target: \texttt{CLICK\ {[}59393{]}}
\item
  offline-side prediction: \texttt{CLICK\ {[}59393{]}}
\item
  online-side prediction: \texttt{CLICK\ {[}63241{]}}
\end{itemize}

So the direct pair is internally consistent with the aggregate
decomposition above:

\begin{itemize}
\tightlist
\item
  \texttt{online\_wf} helps on local value formatting
\item
  \texttt{offline\_wf} is slightly steadier on result-page clicking
\end{itemize}

\subsubsection{Workflow-Text
Difference}

The workflow texts show a structural contrast.

\paragraph{Offline workflow text}

The audited \texttt{kayak} offline workflow contains broader reusable
search primitives:

\begin{itemize}
\tightlist
\item
  enter location for car or hotel search
\item
  select travel dates
\item
  select one-way flight
\item
  search vacation packages
\item
  search flights
\item
  select flight filters
\item
  select hotel filters
\item
  view and select deals
\end{itemize}

These are relatively general travel-site routines. They are not tied to
one exact successful task sequence.

\paragraph{Online workflow text}

The audited \texttt{kayak} online workflow is narrower and more
trajectory-shaped:

\begin{itemize}
\tightlist
\item
  \texttt{search\_for\_cars\_kayak}
\item
  \texttt{select\_location\_kayak}
\item
  \texttt{select\_dates\_for\_rental\_kayak}
\item
  \texttt{search\_flights\_kayak}
\item
  \texttt{browse\_car\_type\_kayak}
\end{itemize}

Compared with the offline library, the online library is:

\begin{itemize}
\tightlist
\item
  more site-specific
\item
  more task-family-specific
\item
  more tightly coupled to the successful test-time routines that
  generated it
\end{itemize}

This supports the interpretation that \texttt{online\_wf} is closer to
the current test distribution, but also less robust as a general
reusable library.

\subsubsection{\texorpdfstring{A7.3 First-Run Target-Site
Underperformance for \texttt{online\_wf}}{A7.3 First-Run Target-Site Underperformance for online\_wf}}

The same logged setting that helps on \texttt{kayak} does not carry over
as a positive first-run result on the two target sites available here.

\paragraph{Tripadvisor}

From the audited paired-case summary:

\begin{itemize}
\tightlist
\item
  \texttt{positive\ =\ 4}
\item
  \texttt{negative\ =\ 18}
\item
  \texttt{net\ =\ -14}
\end{itemize}

From the manually verified tripadvisor negative cases:

\begin{itemize}
\tightlist
\item
  multiple step-0 failures convert a required \texttt{CLICK} into an
  early \texttt{TYPE}
\item
  example tasks that should begin by entering a category or opening
  navigation are pushed into location search instead
\item
  example wrong values include \texttt{Eiffel\ Tower\ Paris},
  \texttt{Jaipur}, \texttt{San\ Francisco}, and \texttt{Egypt}
\end{itemize}

This is consistent with a recorded prompt-level mismatch pattern, but
not by itself a controlled causal proof.

\paragraph{Reddit}

From the audited paired-case summary:

\begin{itemize}
\tightlist
\item
  \texttt{positive\ =\ 2}
\item
  \texttt{negative\ =\ 5}
\item
  \texttt{net\ =\ -3}
\end{itemize}

From the manually verified reddit negative cases:

\begin{itemize}
\tightlist
\item
  the failures are mostly \texttt{CLICK} substitutions
\item
  the induced workflow assumes a small menu of Reddit-specific routines
  such as search, join, sort, and time filter
\item
  these routines are not broadly transferable once the target task
  diverges from the learned community/search template
\end{itemize}

This is weaker than the Tripadvisor case and should be read only as
same-direction exploratory evidence.

\subsubsection{Exploratory
reading}

The evidence supports only a cautious exploratory reading.

\paragraph{Offline workflow}

\begin{itemize}
\tightlist
\item
  derived from training data rather than from the
  model\textquotesingle s own recent test-time trajectories
\item
  broader and more reusable at the text level
\item
  steadier on \texttt{CLICK} grounding and later-stage navigation on
  \texttt{kayak}
\item
  potentially vulnerable when the learned library is semantically
  misaligned with the evaluated task/site
\end{itemize}

\paragraph{Online workflow}

\begin{itemize}
\tightlist
\item
  induced from same-run trajectories in the current implementation, so
  it is closer to the immediate test distribution
\item
  stronger at local TYPE/value guidance on \texttt{kayak}
\item
  narrower and more routine-shaped
\item
  more likely to encode local biases or wrong subroutines that later
  appear in prompt-level mismatch cases
\end{itemize}

\subsubsection{Interpretive summary}

Taken together, the evidence suggests a trade-off rather than a simple
ranking: \texttt{offline\_wf} is broader and steadier on CLICK-side
grounding on \texttt{kayak}, while \texttt{online\_wf} is more
test-proximate and stronger on local TYPE/value guidance there. The
additional target-site readings should be treated as exploratory
context, not as a standalone causal mechanism claim.

\subsubsection{Interpretation limits}

This is still first-run evidence, not a repeated-estimate claim.

In particular:

\begin{itemize}
\tightlist
\item
  the direct \texttt{offline\_wf} vs \texttt{online\_wf} comparison here
  is only on \texttt{kayak}
\item
  the target-site underperformance evidence comes from \texttt{tripadvisor} and
  \texttt{reddit}
\item
  the mechanism reading remains weaker than controlled causal
  identification and should be treated as exploratory
\end{itemize}


% A8
\clearpage
\subsection{Online Small-Data Efficiency}\label{app:A8}

\subsubsection{Overview}

This appendix is retained as an exploratory prefix-level rereading of
logged online results. It should not be used as a primary appendix for
a standalone mechanism claim.

The paper\textquotesingle s online-memory narrative implies that a small
number of trajectories may already be enough to make the induced
workflow useful.

This appendix checks that implication without adding any new experiment.
It asks:

\begin{itemize}
\tightlist
\item
  when \texttt{induce\_steps=1}, does \texttt{online\_wf} show
  measurable gains after only a few induced examples?
\item
  is that early-gain pattern stable across source and target settings?
\end{itemize}

This appendix is derived from the reconstructed online prefix curves,
cross-checked against the audited step-level and paired-case summaries.

\subsubsection{Method}

For each online setting, define:

\begin{itemize}
\tightlist
\item
  \texttt{budget\ =\ i}
\item
  where \texttt{i} is the number of prior examples already used for
  workflow induction when evaluating sample \texttt{i}
\end{itemize}

Under \texttt{induce\_steps=1}:

\begin{itemize}
\tightlist
\item
  \texttt{budget=0} means the very first sample, still without induced
  workflow
\item
  \texttt{budget=1} means one prior example has already been folded
  into the workflow
\item
  \texttt{budget=2} means two prior examples have already been folded in
\end{itemize}

For each site, we compare:

\begin{itemize}
\tightlist
\item
  \texttt{online\_wf}
\item
  against the matching \texttt{no\_workflow}
\end{itemize}

and compute cumulative prefix deltas for:

\begin{itemize}
\tightlist
\item
  \texttt{Elem\ Acc}
\item
  \texttt{Action\ F1}
\item
  \texttt{Step\ SR}
\end{itemize}

\subsubsection{Main result}

\paragraph{\texorpdfstring{Prefix-level \texttt{Step\ SR}
deltas}{Prefix-level Step SR deltas}}

{\def\LTcaptype{} % do not increment counter
\begin{longtable}[]{@{}>{\raggedright\arraybackslash}p{0.20\linewidth}>{\raggedright\arraybackslash}p{0.30\linewidth}>{\raggedright\arraybackslash}p{0.22\linewidth}>{\raggedright\arraybackslash}p{0.24\linewidth}@{}}
\toprule\noalign{}
site & earliest positive prefix & best prefix & final prefix \\
\midrule\noalign{}
\endhead
\bottomrule\noalign{}
\endlastfoot
kayak & \texttt{budget=1}, \texttt{+5.00pp} & \texttt{budget=5},
\texttt{+5.63pp} & \texttt{+5.63pp} \\
tripadvisor & none & best observed still negative (\texttt{budget=1},
\texttt{-3.57pp}) & \texttt{-11.74pp} \\
reddit & \texttt{budget=10}, \texttt{+2.02pp} & \texttt{budget=10},
\texttt{+2.02pp} & \texttt{-2.16pp} \\
\end{longtable}
}

The shortest reading of the prefix table is:

\begin{itemize}
\tightlist
\item
  \textbf{yes on \texttt{kayak}}
\item
  \textbf{no on the harder target-site first run (\texttt{tripadvisor})}
\item
  \textbf{only weakly and temporarily on \texttt{reddit}}
\end{itemize}

So the current evidence does not support a benchmark-wide small-data
claim. It supports a narrower descriptive reading:
\textbf{very-small-budget gains appear on \texttt{kayak}, but the same
pattern is not reproduced on the two available target-site first
runs.}

\subsubsection{Site-by-Site Reading}

\paragraph{\texorpdfstring{\texttt{kayak}}{kayak}}

From the audited online prefix curve:

\begin{itemize}
\tightlist
\item
  after only \textbf{one} induced example, cumulative \texttt{Step\ SR}
  is already \texttt{+5.00pp}
\item
  the gain remains positive throughout the available prefix
\item
  the final prefix reaches \texttt{+5.63pp}
\end{itemize}

This is real early gain, not just a late-stage effect.

Combined with Appendix~\ref{app:A7},
a cautious reading is:

\begin{itemize}
\tightlist
\item
  \texttt{online\_wf} quickly learns a locally useful routine
\item
  on \texttt{kayak}, that routine is close enough to the test
  distribution to help early
\end{itemize}

\paragraph{\texorpdfstring{\texttt{tripadvisor}}{tripadvisor}}

\begin{itemize}
\tightlist
\item
  the first induced example already yields cumulative
  \texttt{Step\ SR\ =\ -3.57pp}
\item
  no positive prefix appears at any small budget
\item
  the full prefix ends at \texttt{-11.74pp}
\end{itemize}

This is the opposite of a small-data efficiency story.

One exploratory explanation is not simply "too little data", but
\textbf{wrong early pattern acquisition}:

\begin{itemize}
\tightlist
\item
  a few early induced trajectories are enough to create a workflow
\item
  but that workflow is already semantically mismatched to the target
  navigation structure
\item
  so small data does not produce early benefit; it produces early bias
\end{itemize}

\paragraph{\texorpdfstring{\texttt{reddit}}{reddit}}

\begin{itemize}
\tightlist
\item
  there is no early gain at budgets \texttt{1,\ 2,\ 3,\ 5}
\item
  the first positive cumulative \texttt{Step\ SR} appears only at
  \texttt{budget=10}
\item
  that gain is small (\texttt{+2.02pp}) and disappears by the final
  prefix
\end{itemize}

So \texttt{reddit} is not a clean "small-data works" case either.

The more cautious interpretation is:

\begin{itemize}
\tightlist
\item
  online induction may need more than a handful of examples before any
  net benefit appears
\item
  even then, the benefit may be unstable
\end{itemize}

\subsubsection{Mechanism
reading}

This pattern is consistent with a broader exploratory reading developed
elsewhere in the project:

\begin{itemize}
\tightlist
\item
  on \texttt{kayak}, \texttt{online\_wf} can acquire a useful local
  routine quickly
\item
  on the two target-site first runs here, the same small-budget pattern
  is not stable
\end{itemize}

So the small-data question should be answered conditionally:

Online workflow induction can show early gains under very small budgets,
but the current first-run evidence does not show the same pattern on
\texttt{tripadvisor} or \texttt{reddit}.

\subsubsection{Interpretive summary}

Under the current Mind2Web first-run evidence, online memory shows
genuine early gains on \texttt{kayak} after only one induced example,
but this very-small-budget benefit is not reproduced on the available
target-site first runs \texttt{tripadvisor} and \texttt{reddit}.

\subsubsection{Interpretation limits}

\begin{itemize}
\tightlist
\item
  This is still a first-run analysis.
\item
  The prefix curve is reconstructed from logged online results, not from
  repeated random seeds.
\item
  The comparison is strongest as a descriptive small-budget pattern, not
  as a formal learning-curve estimate.
\end{itemize}


% B1
\clearpage
\subsection{Positive Case: Kayak Date-Selection Correction}\label{app:B1}

This appendix provides a clean positive case where the workflow
corrects a late-stage click on a matched site. The case was manually
verified against the paired baseline/workflow trajectory logs and the
audited case notes.

\subsubsection{Case metadata}

{\def\LTcaptype{} % do not increment counter
\begin{longtable}[]{@{}>{\raggedright\arraybackslash}p{0.20\linewidth}>{\raggedright\arraybackslash}p{0.76\linewidth}@{}}
\toprule\noalign{}
Field & Value \\
\midrule\noalign{}
\endhead
\bottomrule\noalign{}
\endlastfoot
Website & \texttt{kayak} \\
Setting & \texttt{gpt-4o}, \texttt{test\_task}, \texttt{offline\_wf} vs
\texttt{no\_workflow} \\
Task id & \texttt{4} \\
Raw step index & \texttt{6} \\
Case-study step label & \texttt{8\ /\ 14} \\
Task &
\texttt{check\ cheap\ flights\ from\ NYC\ to\ London\ on\ 23rd\ of\ April\ for\ students\ over\ 18\ years.} \\
\end{longtable}
}

Note: the case-study file uses a different step numbering convention
(\texttt{step=8/14}), while the raw JSON pair index corresponds to raw
step index \texttt{6}.\\
This appendix uses the raw JSON indexing to avoid ambiguity.

\subsubsection{Claim supported}

This case supports a narrow prompt-level claim: under the recorded
\texttt{offline\_wf} condition, workflow memory was already present in
the prompt before raw step \texttt{6}, and the workflow condition
selected the correct date element where the baseline selected a nearby
wrong element.

This case does not by itself isolate workflow insertion as the only
changed factor between control and treatment, so it should be read as
mechanism-consistent evidence rather than strict causal identification.

\subsubsection{What is verified here}

\begin{enumerate}
\def\labelenumi{\arabic{enumi}.}
\tightlist
\item the case belongs to the official \texttt{gpt-4o / test\_task / kayak / offline\_wf} comparison setting;
\item the \texttt{offline\_wf} prompt already contains workflow memory at raw step \texttt{6};
\item the baseline predicts the wrong date element, while the workflow condition predicts the correct one;
\item the corrected action is locally consistent with the date-selection routine visible in the injected workflow text.
\end{enumerate}

\subsubsection{Target and
predictions}

{\def\LTcaptype{} % do not increment counter
\begin{longtable}[]{@{}>{\raggedright\arraybackslash}p{0.28\linewidth}>{\raggedright\arraybackslash}p{0.46\linewidth}>{\raggedright\arraybackslash}p{0.22\linewidth}@{}}
\toprule\noalign{}
Condition & Prediction & Correct? \\
\midrule\noalign{}
\endhead
\bottomrule\noalign{}
\endlastfoot
Target & \texttt{CLICK\ {[}59393{]}} & --- \\
Baseline (\texttt{no\_workflow}) & \texttt{CLICK\ {[}63241{]}} & No \\
Workflow (\texttt{offline\_wf}) & \texttt{CLICK\ {[}59393{]}} & Yes \\
\end{longtable}
}

\subsubsection{Observation excerpt from raw
JSON}

\begin{verbatim}
Observation: `<html> <div menu> <div tab start date calendar input use> ... <div id=63241 button wednesday may 3, 2023> 3 </div> ... <div id=59393 button wednesday may 3, 2023> 3 </div> ... </div> </div> </html>`
\end{verbatim}

\subsubsection{Workflow excerpt present in the
prompt}

\begin{verbatim}
## Workflow 2: Select Travel Dates
Given that you are on the search page for flights, cars, or hotels, this workflow selects the travel dates.
- [generic/button] Start date/End date -> CLICK
- [div/gridcell] {start-date} -> CLICK
- [div/gridcell] {end-date} -> CLICK
\end{verbatim}

\subsubsection{Raw action outputs}

\paragraph{Baseline}

\begin{verbatim}
Action: `CLICK [63241]` ([div] 3 -> CLICK)
\end{verbatim}

\paragraph{Workflow condition}

\begin{verbatim}
Action: `CLICK [59393]` ([button]  May 3, 2023 -> CLICK)
\end{verbatim}

\subsubsection{Interpretation}

Here the baseline selects a nearby distractor with the same visible date
text, while the workflow-conditioned run selects the correct calendar
target. The case therefore illustrates how workflow guidance can
stabilize late-stage element choice when multiple visually similar date
options are present.

\subsubsection{Limitation}

This appendix supports step-level element grounding improvement on a
matched site. It should not be read as evidence that all \texttt{kayak}
gains come from the same workflow fragment.


% B2
\clearpage
\subsection{Positive Case: United Action-Mode Redirection}\label{app:B2}

This appendix provides a clean positive case where workflow changes the
action mode from an incorrect \texttt{CLICK} to a correct
\texttt{TYPE}. The case was manually verified against the paired
baseline/workflow trajectory logs and the audited case notes.

\subsubsection{Case metadata}

{\def\LTcaptype{} % do not increment counter
\begin{longtable}[]{@{}>{\raggedright\arraybackslash}p{0.20\linewidth}>{\raggedright\arraybackslash}p{0.76\linewidth}@{}}
\toprule\noalign{}
Field & Value \\
\midrule\noalign{}
\endhead
\bottomrule\noalign{}
\endlastfoot
Website & \texttt{united} \\
Setting & \texttt{qwen/qwen3.5-397b-a17b}, \texttt{test\_task},
\texttt{lm\_wf} vs \texttt{no\_workflow} \\
Task id & \texttt{3} \\
Raw step index & \texttt{3} \\
Case-study label & \texttt{5\ /\ 27} \\
Task &
\texttt{Find\ a\ basic\ economy\ flight\ +\ hotel\ for\ an\ award\ travel\ from\ las\ vegas\ to\ san\ francisco\ leaving\ and\ returning\ on\ any\ date\ on\ april\ for\ 1\ traveler\ and\ one\ room} \\
\end{longtable}
}

\subsubsection{Claim supported}

This case supports a narrow prompt-level claim: under the recorded
\texttt{lm\_wf} condition, workflow memory was already present in the
prompt before raw step \texttt{3}, and the workflow condition produced
the correct \texttt{TYPE} where the baseline produced an incorrect
\texttt{CLICK}.

This case does not by itself establish workflow insertion as the sole
causal difference between the two conditions, so it should be read as
mechanism-consistent evidence rather than strict causal identification.

\subsubsection{What is verified here}

\begin{enumerate}
\def\labelenumi{\arabic{enumi}.}
\tightlist
\item the case belongs to the official \texttt{qwen/qwen3.5-397b-a17b / test\_task / united / lm\_wf} comparison setting;
\item the \texttt{lm\_wf} prompt already contains workflow memory at raw step \texttt{3};
\item the baseline predicts an incorrect \texttt{CLICK}, while the workflow condition predicts the correct \texttt{TYPE};
\item the corrected \texttt{TYPE} is locally consistent with the location-entry routine visible in the injected workflow text.
\end{enumerate}

\subsubsection{Target and
predictions}

{\def\LTcaptype{} % do not increment counter
\begin{longtable}[]{@{}>{\raggedright\arraybackslash}p{0.28\linewidth}>{\raggedright\arraybackslash}p{0.46\linewidth}>{\raggedright\arraybackslash}p{0.22\linewidth}@{}}
\toprule\noalign{}
Condition & Prediction & Correct? \\
\midrule\noalign{}
\endhead
\bottomrule\noalign{}
\endlastfoot
Target & \texttt{TYPE\ {[}10892{]}\ {[}las\ vegas{]}} & --- \\
Baseline (\texttt{no\_workflow}) & \texttt{CLICK\ {[}10978{]}} & No \\
Workflow (\texttt{lm\_wf}) &
\texttt{TYPE\ {[}10892{]}\ {[}las\ vegas{]}} & Yes \\
\end{longtable}
}

\subsubsection{Observation excerpt from raw
JSON}

\begin{verbatim}
Observation: `<html> <div> <div radiogroup package-type> <label id=10978> <div> <input radio hotel + flight /> <svg img /> </div> Hotel + Flight </label> <label> Hotel + Flight + Car </label> <label> Flight + Car </label> </div> <div> <svg img /> <input id=10892 combobox text where to? /> </div> <div> <div> <svg id=11159 img /> <input id=11161 departing - returning /> </div> <button id=11164 button select traveler information> <div> <svg img /> 2 Adults, 1 Room </div> </button> </div> </div> </html>`
\end{verbatim}

\subsubsection{Workflow excerpt present in the
prompt}

\begin{verbatim}
## enter_flight_locations
Given that you are on the United Airlines booking or search page, this workflow enters the departure and destination city/airport for your flight.
[combobox]  Enter your departing city, airport name, or airpor... -> TYPE: {origin-city}
[button]  {origin-option} -> CLICK
[combobox]  Enter your destination city, airport name, or airp... -> TYPE: {destination-city}
[button]  {destination-option} -> CLICK
\end{verbatim}

\subsubsection{Raw action outputs}

\paragraph{Baseline}

\begin{verbatim}
Action: `CLICK [10978]` ([label]  Hotel + Flight -> CLICK)
\end{verbatim}

\paragraph{Workflow condition}

\begin{verbatim}
Action: `TYPE [10892] [las vegas]` ([combobox]  Where to? -> TYPE: las vegas)
\end{verbatim}

\subsubsection{Interpretation}

Here the baseline clicks a package-type control, while the
workflow-conditioned run types into the correct origin-city combobox.
The case therefore illustrates action-mode redirection: the workflow is
useful not because it changes a value, but because it steers the model
toward the right kind of action.

\subsubsection{Limitation}

This is evidence for step-level action-mode correction. It should not be
used to claim that all positive cases are driven by the same mechanism.


% B4
\clearpage
\subsection{Negative Case: Sixflags Workflow-Content Mismatch}\label{app:B4}

This appendix provides a clean negative case where a semantically
mismatched workflow harms performance. The case was manually verified
against the paired baseline/workflow trajectory logs and the audited
case notes.

\subsubsection{Case metadata}

{\def\LTcaptype{} % do not increment counter
\begin{longtable}[]{@{}>{\raggedright\arraybackslash}p{0.20\linewidth}>{\raggedright\arraybackslash}p{0.76\linewidth}@{}}
\toprule\noalign{}
Field & Value \\
\midrule\noalign{}
\endhead
\bottomrule\noalign{}
\endlastfoot
Website & \texttt{sixflags} \\
Setting & \texttt{gpt-4o}, \texttt{test\_task}, \texttt{offline\_wf} vs
\texttt{no\_workflow} \\
Result-file task id & \texttt{2} \\
Raw step index & \texttt{0} \\
Case-study label & \texttt{0\ /\ 6} \\
Current step task &
\texttt{Apply\ for\ a\ job\ on\ the\ Six\ Flags\ White\ Water\ park} \\
\end{longtable}
}

\subsubsection{Claim supported}

This case supports a narrow prompt-level claim: under the recorded
\texttt{offline\_wf} condition, workflow memory was already present in
the prompt before raw step \texttt{0}, and the workflow condition
produced the wrong park-selection click where the baseline produced the
correct \texttt{Jobs} click.

This case does not by itself isolate workflow insertion as the only
changed factor between treatment and control, so it should be read as
mechanism-consistent evidence rather than strict causal identification.

\subsubsection{What is verified here}

\begin{enumerate}
\def\labelenumi{\arabic{enumi}.}
\tightlist
\item the case belongs to the audited \texttt{gpt-4o / test\_task / sixflags / offline\_wf} setting;
\item the current step in the raw JSON prompt is a jobs-related task, not the earlier balance-sheet task that appears in the file-level case-study index;
\item the baseline predicts the correct \texttt{Jobs} navigation click, while the workflow condition predicts the wrong \texttt{Browse the Parks Below} click;
\item the wrong click is locally consistent with the park-selection workflows present in the prompt.
\end{enumerate}

\subsubsection{Target and
predictions}

{\def\LTcaptype{} % do not increment counter
\begin{longtable}[]{@{}>{\raggedright\arraybackslash}p{0.28\linewidth}>{\raggedright\arraybackslash}p{0.46\linewidth}>{\raggedright\arraybackslash}p{0.22\linewidth}@{}}
\toprule\noalign{}
Condition & Prediction & Correct? \\
\midrule\noalign{}
\endhead
\bottomrule\noalign{}
\endlastfoot
Target & \texttt{CLICK\ {[}103{]}} & --- \\
Baseline (\texttt{no\_workflow}) & \texttt{CLICK\ {[}103{]}} & Yes \\
Workflow (\texttt{offline\_wf}) & \texttt{CLICK\ {[}1042{]}} & No \\
\end{longtable}
}

\subsubsection{Observation excerpt from raw
JSON}

\begin{verbatim}
Task: Apply for a job on the Six Flags White Water park
Trajectory:
Observation: `<html> <div> <div> <nav navigation> <a id=103> Jobs </a> </nav> <nav navigation> <a id=766> Jobs </a> </nav> </div> <div> <div select> <button id=1042> <div> Browse the Parks Below </div> </button> <button id=1045> Go! </button> </div> <form news letters form> <div> <label> First Name </label> <input id=2747 text andy /> </div> </form> </div> </div> </html>`
\end{verbatim}

\subsubsection{Workflow excerpt present in the
prompt}

\begin{verbatim}
## select_park
Given that you are on the Six Flags website, this workflow selects a specific park from a list of options.
[button]  Browse the Parks Below -> CLICK
[span]  {park-name} -> CLICK
[button]  Go! -> CLICK

## select_event_options
Given that you are on a selected Six Flags park's page, this workflow navigates to see all events available at that park.
[link]  Events [menu] -> CLICK
[link]  See all events [next] -> CLICK
\end{verbatim}

\subsubsection{Raw action outputs}

\paragraph{Baseline}

\begin{verbatim}
Action: `CLICK [103]`
\end{verbatim}

\paragraph{Workflow condition}

\begin{verbatim}
Action: `CLICK [1042]`
\end{verbatim}

\subsubsection{Interpretation}

Here the baseline takes the correct \texttt{Jobs} navigation link, while
the workflow-conditioned run follows a park-selection pattern
(\texttt{Browse the Parks Below}) that belongs to a different task
family. The case therefore illustrates workflow-content mismatch rather
than a generic click error.

\subsubsection{Limitation}

This example is best used as evidence for workflow-content mismatch. It
should not be overstated as proof that the model always follows workflow
over local reasoning.


% B5
\clearpage
\subsection{Negative Case: Sixflags Template-Induced Step Skipping}\label{app:B5}

This appendix provides a negative case where the workflow appears to
encourage a premature \texttt{Book\ Now} click instead of the correct
date-selection step. The case was manually verified against the paired
baseline/workflow trajectory logs and the audited case notes.

\subsubsection{Case metadata}

{\def\LTcaptype{} % do not increment counter
\begin{longtable}[]{@{}>{\raggedright\arraybackslash}p{0.20\linewidth}>{\raggedright\arraybackslash}p{0.76\linewidth}@{}}
\toprule\noalign{}
Field & Value \\
\midrule\noalign{}
\endhead
\bottomrule\noalign{}
\endlastfoot
Website & \texttt{sixflags} \\
Setting & \texttt{gpt-4o}, \texttt{test\_task}, \texttt{offline\_wf} vs
\texttt{no\_workflow} \\
Task id & \texttt{1} \\
Raw step index & \texttt{5} \\
Case-study label & \texttt{Case\ 2:\ task=1,\ step=5/14} \\
Current step task &
\texttt{Find\ a\ place\ to\ stay\ near\ Great\ Escape\ New\ Park\ from\ April\ 21\ to\ April\ 24\ for\ 2\ adults\ and\ 1\ kid,\ and\ book\ the\ cheapest\ themed\ room.} \\
\end{longtable}
}

\subsubsection{Claim supported}

This case supports a narrow prompt-level claim: under the recorded
\texttt{offline\_wf} condition, workflow memory was already present in
the prompt before raw step \texttt{5}, and the workflow condition
produced a premature purchase-oriented \texttt{CLICK} where the
baseline took the correct date-selection action.

This case does not isolate workflow insertion as the only changed factor
between treatment and control, so it should be read as
mechanism-consistent evidence rather than strict causal identification.

\subsubsection{What is verified here}

\begin{enumerate}
\def\labelenumi{\arabic{enumi}.}
\tightlist
\item the case belongs to the official \texttt{gpt-4o / test\_task / sixflags / offline\_wf} comparison setting;
\item the audited raw step is part of the Great Escape lodging task recorded in both paired result files;
\item the \texttt{offline\_wf} prompt already contains workflow memory at raw step \texttt{5};
\item the workflow condition clicks \texttt{Book Now} too early, while the baseline correctly clicks the arrival-date input, and this premature click is locally consistent with the purchase-oriented routine visible in the injected workflow text.
\end{enumerate}

\subsubsection{Target and
predictions}

{\def\LTcaptype{} % do not increment counter
\begin{longtable}[]{@{}>{\raggedright\arraybackslash}p{0.28\linewidth}>{\raggedright\arraybackslash}p{0.46\linewidth}>{\raggedright\arraybackslash}p{0.22\linewidth}@{}}
\toprule\noalign{}
Condition & Prediction & Correct? \\
\midrule\noalign{}
\endhead
\bottomrule\noalign{}
\endlastfoot
Target & \texttt{CLICK\ {[}46822{]}} & --- \\
Baseline (\texttt{no\_workflow}) & \texttt{CLICK\ {[}46822{]}} & Yes \\
Workflow (\texttt{offline\_wf}) & \texttt{CLICK\ {[}47084{]}} & No \\
\end{longtable}
}

\subsubsection{Observation excerpt from raw
JSON}

\begin{verbatim}
Observation: `<html> <div> <a id=47084 book now> Book Now </a> <div> <div> <label> Check in date </label> <input id=46822 text startdate select arrival date /> </div> <div> <label> Check out date </label> <input id=47302 text enddate select departure date /> </div> <input hotelid 6491 /> <button id=47310 submit> Check Rates </button> </div> ... </html>`
\end{verbatim}

\subsubsection{Workflow excerpt present in the
prompt}

\begin{verbatim}
## buy_group_tickets
Given that you are on a selected Six Flags park's page, this workflow buys group tickets by entering necessary group information.
[link]  Groups -> CLICK
[button]  Buy Now -> CLICK
[button]  {date} -> CLICK
[button]  Buy Now -> CLICK
[input]   -> TYPE: {group-name}
...
\end{verbatim}

\subsubsection{Raw action outputs}

\paragraph{Baseline}

\begin{verbatim}
Action: `CLICK [46822]` (select Check in date -> CLICK)
\end{verbatim}

\paragraph{Workflow condition}

\begin{verbatim}
Action: `CLICK [47084]` ([link]  Book Now -> CLICK)
\end{verbatim}

\subsubsection{Interpretation}

Here the baseline selects the required check-in date field, while the
workflow-conditioned run jumps directly to \texttt{Book Now}. The case
therefore illustrates template-induced step skipping: a compressed
purchase routine pulls the model to a later action before a necessary
local substep has been completed.

\subsubsection{Limitation}

This appendix supports a ``template-induced step skipping''
interpretation for this specific case. It should not be generalized into
a universal rule without additional audited examples.


% B6
\clearpage
\subsection{Negative Case: Tripadvisor Online-Workflow Mismatch}\label{app:B6}

This appendix documents a prompt-level mismatch case under the
\texttt{test\_website / tripadvisor / online\_wf} setting. The case was
manually verified against the paired baseline/workflow trajectory logs
and the audited case notes.

\subsubsection{Case metadata}

{\def\LTcaptype{} % do not increment counter
\begin{longtable}[]{@{}>{\raggedright\arraybackslash}p{0.20\linewidth}>{\raggedright\arraybackslash}p{0.76\linewidth}@{}}
\toprule\noalign{}
Field & Value \\
\midrule\noalign{}
\endhead
\bottomrule\noalign{}
\endlastfoot
Website & \texttt{tripadvisor} \\
Setting & \texttt{qwen/qwen3.5-397b-a17b}, \texttt{test\_website},
\texttt{online\_wf} vs \texttt{no\_workflow} \\
Task id & \texttt{2} \\
Raw step index & \texttt{0} \\
Case-study label & \texttt{0\ /\ 13} \\
Task &
\texttt{Find\ a\ vacation\ rental\ with\ 2\ bedrooms\ and\ 2\ bathrooms\ for\ the\ 15\ to\ 20\ June\ period,\ in\ Kanha\ National\ Park,\ India,\ and\ book\ the\ cheapest\ one.} \\
\end{longtable}
}

\subsubsection{Claim supported}

This case supports a narrow prompt-level claim: under the recorded
\texttt{online\_wf} condition, workflow memory was already present in
the prompt before the first action of task \texttt{2}, and the
treatment condition produced an incorrect early \texttt{TYPE} where the
baseline produced the correct \texttt{CLICK}.

This case does not by itself isolate workflow insertion as the only
changed factor between treatment and control, so it should be read as
mechanism-consistent evidence rather than strict causal identification.

\subsubsection{What is verified here}

\begin{enumerate}
\def\labelenumi{\arabic{enumi}.}
\tightlist
\item the case belongs to the official \texttt{C2 / test\_website / tripadvisor / online\_wf} run setting;
\item the \texttt{online\_wf} prompt already contains workflow memory at task \texttt{2}, step \texttt{0};
\item the baseline predicts the correct first-step \texttt{CLICK}, while the \texttt{online\_wf} condition predicts an incorrect \texttt{TYPE};
\item the wrong \texttt{TYPE} is locally consistent with the search-first routine visible in the injected workflow message.
\end{enumerate}

\subsubsection{Target and
predictions}

{\def\LTcaptype{} % do not increment counter
\begin{longtable}[]{@{}>{\raggedright\arraybackslash}p{0.28\linewidth}>{\raggedright\arraybackslash}p{0.46\linewidth}>{\raggedright\arraybackslash}p{0.22\linewidth}@{}}
\toprule\noalign{}
Condition & Prediction & Correct? \\
\midrule\noalign{}
\endhead
\bottomrule\noalign{}
\endlastfoot
Target & \texttt{CLICK\ {[}631{]}} & --- \\
Baseline (\texttt{no\_workflow}) & \texttt{CLICK\ {[}631{]}} & Yes \\
Workflow (\texttt{online\_wf}) &
\texttt{TYPE\ {[}808{]}\ {[}Kanha\ National\ Park,\ India{]}} & No \\
\end{longtable}
}

\subsubsection{Observation excerpt from raw
JSON}

\begin{verbatim}
Task: Find a vacation rental with 2 bedrooms and 2 bathrooms for the 15 to 20 June period, in Kanha National Park, India, and book the cheapest one.
Trajectory:
Observation: `<html> <div> <form search> ... <a id=631> <span> Vacation Rentals </span> </a> ... <input id=808 searchbox search q where to? /> ... </html>`
\end{verbatim}

\subsubsection{Workflow message actually present in the
prompt}

\begin{verbatim}
## Summary Workflows:
## search_destination
Given that you are on the TripAdvisor search page, this workflow searches for and selects a specific destination or attraction.
[input]  Search where to? -> TYPE: {destination}
[option]  {destination_suggestion} -> CLICK
\end{verbatim}

\subsubsection{Raw action outputs}

\paragraph{Baseline}

\begin{verbatim}
Action: `CLICK [631]` ([link]  Vacation Rentals -> CLICK)
\end{verbatim}

\paragraph{Workflow condition}

\begin{verbatim}
Action: `TYPE [808] [Kanha National Park, India]`
\end{verbatim}

\subsubsection{Interpretation}

Here the baseline takes the correct \texttt{Vacation Rentals}
navigation click, while the workflow-conditioned run types a location
into the search box. Because the treatment prompt already contains a
search-first workflow routine, the error is consistent with a
workflow-conditioned entry-point mismatch.

\subsubsection{Limitation}

This case does not prove workflow provenance beyond the current
prefix-level online setting, does not show that the workflow was induced
only from successful prior trajectories, and does not isolate workflow
insertion as the only changed prompt factor. It is therefore best used
together with aggregated target-site result tracing, not as a standalone causal
proof.


% C1
\clearpage
\subsection{LM vs Rule Workflow Text Evidence}\label{app:C1}

This appendix supports the claim that LM-induced workflows are more
abstract and reusable than rule-induced workflows at the text level.
The evidence below is drawn from the audited workflow-text comparison.

\subsubsection{Summary table from the source
output}

{\def\LTcaptype{} % do not increment counter
\begin{longtable}[]{@{}>{\raggedright\arraybackslash}p{0.11\linewidth}>{\raggedleft\arraybackslash}p{0.10\linewidth}>{\raggedleft\arraybackslash}p{0.10\linewidth}>{\raggedleft\arraybackslash}p{0.10\linewidth}>{\raggedleft\arraybackslash}p{0.10\linewidth}>{\raggedleft\arraybackslash}p{0.10\linewidth}>{\raggedleft\arraybackslash}p{0.10\linewidth}>{\raggedright\arraybackslash}p{0.23\linewidth}@{}}
\toprule\noalign{}
Website & LM workflows & Rule workflows & LM avg steps/WF & Rule avg
steps/WF & LM placeholders & Rule concrete values & Source verdict \\
\midrule\noalign{}
\endhead
\bottomrule\noalign{}
\endlastfoot
\texttt{kayak} & 8 & 17 & 1.8 & 12.5 & 13 & 25 &
\texttt{SUPPORTED\ (4/4\ indicators)} \\
\texttt{newegg} & 6 & 19 & 1.7 & 7.8 & 6 & 16 &
\texttt{SUPPORTED\ (4/4\ indicators)} \\
\texttt{united} & 6 & 24 & 2.7 & 9.0 & 15 & 41 &
\texttt{SUPPORTED\ (4/4\ indicators)} \\
\end{longtable}
}

\subsubsection{Source excerpts}

\paragraph{Kayak}

\begin{verbatim}
[Abstraction Indicators]
  + LM uses more placeholders (13 vs 0)
  + Rule retains more concrete values (25 vs 0)
  + LM workflows are more compact (1.8 vs 12.5 steps/wf)
  + LM has more abstract/reusable workflows (8 vs 0)

[Paper Claim: 'LM induction produces more abstract, reusable sub-routines']
  -> SUPPORTED (4/4 indicators)
\end{verbatim}

\paragraph{Newegg}

\begin{verbatim}
[Abstraction Indicators]
  + LM uses more placeholders (6 vs 0)
  + Rule retains more concrete values (16 vs 0)
  + LM workflows are more compact (1.7 vs 7.8 steps/wf)
  + LM has more abstract/reusable workflows (6 vs 0)
\end{verbatim}

\paragraph{United}

\begin{verbatim}
[Abstraction Indicators]
  + LM uses more placeholders (15 vs 0)
  + Rule retains more concrete values (41 vs 0)
  + LM workflows are more compact (2.7 vs 9.0 steps/wf)
  + LM has more abstract/reusable workflows (6 vs 0)
\end{verbatim}

\subsubsection{Representative workflow
excerpts}

\paragraph{\texorpdfstring{LM workflow example
(\texttt{united})}{LM workflow example (united)}}

\begin{verbatim}
## enter_flight_locations
[combobox]  Enter your departing city, airport name, or airpor... -> TYPE: {origin-city}
[button]  {origin-option} -> CLICK
[combobox]  Enter your destination city, airport name, or airp... -> TYPE: {destination-city}
[button]  {destination-option} -> CLICK
\end{verbatim}

\paragraph{\texorpdfstring{Rule workflow example
(\texttt{united})}{Rule workflow example (united)}}

\begin{verbatim}
Workflow 5: Concrete Trajectory for Find a round trip from Phoenix to Miami with...
Steps: [combobox]  Flying from -> TYPE: Phoenix; [button]  Phoenix, AZ, US (PHX) -> CLICK; [button]  Search -> CLICK ...
\end{verbatim}

\subsubsection{Interpretation limits}

\begin{itemize}
\tightlist
\item
  These counts currently follow the
  \texttt{wf\_text\_compare\_output.txt} counting rule, which includes
  the \texttt{Summary\ Workflows} block.
\item
  This appendix is text-structural evidence. It should be paired with
  score tables if the main text wants to discuss performance consequences.
\end{itemize}


% C2
\clearpage
\subsection{Concrete-Value Counting Rule}\label{app:C2}

This appendix fixes the counting drift around ``concrete values'' in
\texttt{LM} workflow texts. The rule below is derived from the audited
workflow-text comparison and manual inspection of the \texttt{LM}
workflow libraries.

\subsubsection{Recommended counting
rule}

For the LM workflows, count a token as a \texttt{concrete\ value} only
if it is a task-instantiated or user-instantiated value that could vary
by task, such as:

\begin{itemize}
\tightlist
\item
  city names
\item
  dates
\item
  product names
\item
  confirmation numbers
\item
  personal identifiers
\end{itemize}

Do \textbf{not} count fixed UI labels as concrete task values, such as:

\begin{itemize}
\tightlist
\item
  \texttt{Increment}
\item
  \texttt{View\ Deal}
\item
  \texttt{APPLY}
\item
  \texttt{TRAVEL\ INFO}
\item
  \texttt{Find\ cars\ button.}
\end{itemize}

Under this rule, the LM workflow libraries for \texttt{kayak},
\texttt{newegg}, and \texttt{united} contain:

{\def\LTcaptype{} % do not increment counter
\begin{longtable}[]{@{}>{\raggedright\arraybackslash}p{0.36\linewidth}>{\centering\arraybackslash}p{0.60\linewidth}@{}}
\toprule\noalign{}
Website & LM concrete values \\
\midrule\noalign{}
\endhead
\bottomrule\noalign{}
\endlastfoot
\texttt{kayak} & 0 \\
\texttt{newegg} & 0 \\
\texttt{united} & 0 \\
\end{longtable}
}

\subsubsection{Why this rule is
needed}

The earlier drift comes from mixing two things:

\begin{enumerate}
\def\labelenumi{\arabic{enumi}.}
\tightlist
\item
  \textbf{task-variable values}, which should count as concrete values
\item
  \textbf{fixed UI labels}, which should not
\end{enumerate}

The current LM workflow libraries are parameterized almost entirely with
placeholders such as:

\begin{itemize}
\tightlist
\item
  \texttt{\{location\_name\}}
\item
  \texttt{\{search-term\}}
\item
  \texttt{\{origin-city\}}
\item
  \texttt{\{destination-city\}}
\item
  \texttt{\{start-date\}}
\end{itemize}

\subsubsection{Supporting excerpts}

\paragraph{Kayak}

\begin{verbatim}
[textbox]  {location_field} -> TYPE: {location_name}
[span]  {location_suggestion} -> CLICK
[button]  {date_input} -> CLICK
[div]  {date_day} -> CLICK
\end{verbatim}

\paragraph{Newegg}

\begin{verbatim}
[searchbox]  Search Site -> TYPE: {search-term}
[textbox]  price to -> TYPE: {max-price}
[span]  {filter-option} -> CLICK
[link]  {category-link} -> CLICK
\end{verbatim}

\paragraph{United}

\begin{verbatim}
[combobox]  Enter your departing city, airport name, or airpor... -> TYPE: {origin-city}
[button]  {origin-option} -> CLICK
[combobox]  Enter your destination city, airport name, or airp... -> TYPE: {destination-city}
[button]  {destination-option} -> CLICK
\end{verbatim}

\subsubsection{Boundary note}

If a later section wants to count fixed UI labels separately, it should
use a different name such as \texttt{fixed\ interface\ literals}, not
\texttt{concrete\ values}.


% C3
\clearpage
\subsection{Workflow Counting and Site-Feature Definitions}\label{app:C3}

This appendix fixes the counting drift around LM workflow counts and
average steps per workflow. The definitions below are reconstructed
from the audited workflow-text comparison and manual counting of the
\texttt{LM} workflow libraries.

\subsubsection{Two valid counting
conventions}

There are two defensible workflow-count conventions in the current
materials:

\begin{enumerate}
\def\labelenumi{\arabic{enumi}.}
\item
  \textbf{Include \texttt{Summary\ Workflows} as a workflow block}

  \begin{itemize}
  \tightlist
  \item
    This is the convention implicitly used by
    \texttt{wf\_text\_compare\_output.txt}
  \item
    Counts: \texttt{kayak\ 8}, \texttt{newegg\ 6}, \texttt{united\ 6}
  \end{itemize}
\item
  \textbf{Exclude \texttt{Summary\ Workflows} and count only executable
  named workflows}

  \begin{itemize}
  \tightlist
  \item
    This is the convention used in \texttt{C5}
  \item
    Counts: \texttt{kayak\ 7}, \texttt{newegg\ 5}, \texttt{united\ 5}
  \end{itemize}
\end{enumerate}

\subsubsection{Conversion table}

{\def\LTcaptype{} % do not increment counter
\begin{longtable}[]{@{}>{\raggedright\arraybackslash}p{0.14\linewidth}>{\raggedleft\arraybackslash}p{0.22\linewidth}>{\raggedleft\arraybackslash}p{0.14\linewidth}>{\raggedleft\arraybackslash}p{0.14\linewidth}>{\raggedleft\arraybackslash}p{0.14\linewidth}>{\raggedleft\arraybackslash}p{0.18\linewidth}@{}}
\toprule\noalign{}
Website & Total steps & Count incl. summary & Avg steps/WF incl. summary
& Count excl. summary & Avg steps/WF excl. summary \\
\midrule\noalign{}
\endhead
\bottomrule\noalign{}
\endlastfoot
\texttt{kayak} & 14 & 8 & 1.75 & 7 & 2.00 \\
\texttt{newegg} & 10 & 6 & 1.67 & 5 & 2.00 \\
\texttt{united} & 16 & 6 & 2.67 & 5 & 3.20 \\
\end{longtable}
}

\subsubsection{Supporting excerpts}

\paragraph{Kayak}

\begin{verbatim}
## Summary Workflows:
## enter_search_location
...
## select_calendar_date
...
## apply_amenity_filter
...
## sort_search_results
...
## set_flight_trip_type
...
## adjust_guests_rooms
...
## book_selected_option
...
\end{verbatim}

\paragraph{Newegg}

\begin{verbatim}
Website: Shopping,Digital,newegg
## Summary Workflows:
## search_product
...
## apply_price_filter
...
## apply_selection_filter
...
## navigate_category_links
...
## open_site_menu_and_select
...
\end{verbatim}

\subsubsection{Recommended use}

\begin{itemize}
\tightlist
\item
  Use \textbf{include-summary} counts when comparing directly against
  \texttt{wf\_text\_compare\_output.txt}.
\item
  Use \textbf{exclude-summary} counts when discussing \texttt{C5}
  quality metrics.
\item
  Do not mix the two in one table without labeling the convention.
\end{itemize}


% C4
\clearpage
\subsection{Site-Feature Qualitative Coding}\label{app:C4}

This appendix derives site-feature observations from the audited
materials themselves, rather than reverse-engineering evidence for
prewritten claims.

This appendix is meant to support or challenge the qualitative rows in
\texttt{§10.1} of the main text,
especially:

\begin{itemize}
\tightlist
\item
  \texttt{WF\ specificity}
\item
  \texttt{Task\ diversity}
\end{itemize}

The coding below is based on audited offline workflow texts,
workflow-text summaries, step-level summaries, and manually reviewed
case examples.

Boundary:

\begin{itemize}
\tightlist
\item
  This is a qualitative coding memo, not a hard quantitative appendix.
\item
  The goal is to surface what the audited materials actually show about
  site structure and task-family alignment.
\end{itemize}

\subsubsection{Coding questions}

Instead of starting from the paper's conclusions, this appendix asks
four qualitative questions:

\begin{enumerate}
\def\labelenumi{\arabic{enumi}.}
\item
  \textbf{How specific is the workflow library?}

  \begin{itemize}
  \tightlist
  \item
    Does it mostly encode reusable sub-routines, or end-to-end task
    scripts?
  \item
    Does it rely on placeholders or concrete values?
  \end{itemize}
\item
  \textbf{How broad is the observed task family at evaluation time?}

  \begin{itemize}
  \tightlist
  \item
    Across the case studies and workflow failures/successes, do tasks
    cluster around one operational family, or span unrelated site
    functions?
  \end{itemize}
\item
  \textbf{How aligned are workflow families and observed task families?}

  \begin{itemize}
  \tightlist
  \item
    Even if a workflow library is abstract, does it cover the dominant
    task family seen in evaluation?
  \end{itemize}
\item
  \textbf{What mechanism implication follows?}

  \begin{itemize}
  \tightlist
  \item
    Is AWM likely to help through reusable guidance, fail through
    mismatch, or contribute little because baseline headroom is small?
  \end{itemize}
\end{enumerate}

\subsubsection{Coding dimensions}

The following labels are used as qualitative summaries:

\begin{itemize}
\item
  \textbf{WF specificity}

  \begin{itemize}
  \tightlist
  \item
    \texttt{Low}: mostly short parameterized primitives
  \item
    \texttt{Medium}: mixed primitives plus narrow domain patterns
  \item
    \texttt{High}: long or strongly task-family-specific sequences
  \end{itemize}
\item
  \textbf{Observed task-family breadth}

  \begin{itemize}
  \tightlist
  \item
    \texttt{Low}: tasks mostly stay within one narrow operational family
  \item
    \texttt{Medium}: tasks vary, but still share a dominant site
    interaction pattern
  \item
    \texttt{High}: tasks span multiple unrelated functional areas
  \end{itemize}
\item
  \textbf{WF-task alignment}

  \begin{itemize}
  \tightlist
  \item
    \texttt{Strong}: workflows clearly target the dominant eval task
    family
  \item
    \texttt{Partial}: workflows cover only part of the observed task
    family
  \item
    \texttt{Weak}: workflows and eval task families visibly diverge
  \end{itemize}
\end{itemize}

\subsubsection{Site-by-site source
observations}

\paragraph{Kayak}

\textbf{Workflow-text observations}

\begin{itemize}
\tightlist
\item
  The audited offline workflow library contains short reusable
  primitives:

  \begin{itemize}
  \tightlist
  \item
    \texttt{Enter\ Location\ for\ Car\ or\ Hotel\ Search}
  \item
    \texttt{Select\ Travel\ Dates}
  \item
    \texttt{Select\ One-Way\ Flight}
  \item
    \texttt{Select\ Hotel\ Filters}
  \item
    \texttt{View\ and\ Select\ Deals}
  \end{itemize}
\item
  Most workflows are 1 to 4 steps and heavily parameterized:

  \begin{itemize}
  \tightlist
  \item
    \texttt{\{your-location\}}
  \item
    \texttt{\{start-date\}}
  \item
    \texttt{\{filter\}}
  \item
    \texttt{\{best-popup-option\}}
  \end{itemize}
\item
  This is consistent with the workflow-text audit, where \texttt{kayak}
  \texttt{LM} workflows are short and placeholder-heavy.
\end{itemize}

\textbf{Observed eval-task observations}

\begin{itemize}
\tightlist
\item
  Positive cases cover:

  \begin{itemize}
  \tightlist
  \item
    hotel search with amenity constraints
  \item
    cheap flights from NYC to London
  \item
    Hawaii package selection
  \end{itemize}
\item
  Across these cases, the dominant pattern is still a travel-search
  pipeline:

  \begin{itemize}
  \tightlist
  \item
    enter location
  \item
    pick date
  \item
    filter or sort
  \item
    choose a deal
  \end{itemize}
\end{itemize}

\textbf{Coding}

\begin{itemize}
\tightlist
\item
  WF specificity: \texttt{Low}
\item
  Observed task-family breadth: \texttt{Medium}
\item
  WF-task alignment: \texttt{Strong}
\end{itemize}

\textbf{Mechanism implication}

Kayak looks like the cleanest case where abstract sub-routines line up
with a diverse but still coherent search-oriented task family. This
supports the claim that reusable workflows help when variation happens
inside one stable interaction grammar.

\paragraph{Newegg}

\textbf{Workflow-text observations}

\begin{itemize}
\tightlist
\item
  The audited offline workflow library is still fairly abstract, but
  less purely
  primitive than kayak:

  \begin{itemize}
  \tightlist
  \item
    \texttt{search\_and\_apply\_filters}
  \item
    \texttt{search\_and\_sort\_items}
  \item
    \texttt{shopping\_cart\_management}
  \item
    \texttt{build\_custom\_pc}
  \item
    \texttt{custom\_pc\_finder}
  \end{itemize}
\item
  Some workflows are generic search/filter modules; others are narrower
  feature flows (\texttt{build\_custom\_pc},
  \texttt{custom\_pc\_finder}).
\item
  The workflow-text comparison still shows the LM version as more
  abstract than the rule-induced counterpart, but the workflow family is
  more commerce-specific than kayak's travel primitives.
\end{itemize}

\textbf{Observed eval-task observations}

\begin{itemize}
\tightlist
\item
  Positive cases include:

  \begin{itemize}
  \tightlist
  \item
    projector search
  \item
    gaming desktop search
  \item
    cart-related operations
  \item
    bluetooth mouse sorting
  \end{itemize}
\item
  These tasks vary, but the interaction family remains recognizably
  e-commerce:

  \begin{itemize}
  \tightlist
  \item
    search
  \item
    filter
  \item
    sort
  \item
    cart management
  \end{itemize}
\end{itemize}

\textbf{Coding}

\begin{itemize}
\tightlist
\item
  WF specificity: \texttt{Medium-Low}
\item
  Observed task-family breadth: \texttt{Medium}
\item
  WF-task alignment: \texttt{Strong}
\end{itemize}

\textbf{Mechanism implication}

Newegg still supports reusable workflow guidance, but in a narrower
commerce grammar than kayak. This helps explain why both LM and Rule can
be competitive here: the task family is varied enough to benefit from
abstractions, yet stable enough that some concrete routines still
transfer.

\paragraph{Budget}

\textbf{Workflow-text observations}

\begin{itemize}
\tightlist
\item
  The audited offline workflow library mixes several distinct site
  functions:

  \begin{itemize}
  \tightlist
  \item
    rental location search
  \item
    vehicle/extras selection
  \item
    pickup/return dates
  \item
    reservation lookup
  \item
    careers search
  \item
    deals/offers browsing
  \end{itemize}
\item
  Several workflows are more end-to-end and function-specific than
  kayak/newegg:

  \begin{itemize}
  \tightlist
  \item
    \texttt{search\_for\_job\_in\_usa\_finance}
  \item
    \texttt{find\_and\_view\_reservation}
  \item
    \texttt{learn\_about\_deals\_and\_offers}
  \end{itemize}
\item
  This makes the library not just longer on average, but also more
  functionally fragmented.
\end{itemize}

\textbf{Observed eval-task observations}

\begin{itemize}
\tightlist
\item
  Negative and positive cases together show tasks spanning:

  \begin{itemize}
  \tightlist
  \item
    insurance policy certificates
  \item
    car-for-sale browsing
  \item
    rental booking
  \item
    finance job search
  \item
    road-trip ideas / travel content
  \end{itemize}
\item
  The observed task family is therefore much broader than a single
  rental-booking loop.
\end{itemize}

\textbf{Coding}

\begin{itemize}
\tightlist
\item
  WF specificity: \texttt{High}
\item
  Observed task-family breadth: \texttt{High}
\item
  WF-task alignment: \texttt{Weak}
\end{itemize}

\textbf{Mechanism implication}

Budget is the clearest source-based example where the workflow library
is internally fragmented and the observed eval tasks span unrelated
functional areas. This supports a mismatch interpretation more strongly
than a simple ``workflow count too small'' story.

\paragraph{Sixflags}

\textbf{Workflow-text observations}

\begin{itemize}
\tightlist
\item
  The audited offline workflow library is short and park-centric:

  \begin{itemize}
  \tightlist
  \item
    \texttt{select\_park}
  \item
    \texttt{browse\_ticket\_options}
  \item
    \texttt{buy\_group\_tickets}
  \item
    \texttt{select\_event\_options}
  \item
    \texttt{find\_park\_information}
  \end{itemize}
\item
  These workflows are not highly specific in the same sense as budget;
  instead, they are shallow templates around park navigation and
  purchase flows.
\item
  The file suggests a narrow operational grammar rather than a broad
  feature library.
\end{itemize}

\textbf{Observed eval-task observations}

\begin{itemize}
\tightlist
\item
  Case studies show tasks involving:

  \begin{itemize}
  \tightlist
  \item
    park security policy
  \item
    buying a single day pass
  \item
    financial statements
  \item
    park hours
  \item
    VIP tour confirmation
  \end{itemize}
\item
  So the observed eval task family is not simply ``park and ticket
  variants only.'' It includes both on-family park-navigation tasks and
  out-of-family informational/financial tasks.
\end{itemize}

\textbf{Coding}

\begin{itemize}
\tightlist
\item
  WF specificity: \texttt{Medium}
\item
  Observed task-family breadth: \texttt{Medium}
\item
  WF-task alignment: \texttt{Partial\ to\ Weak}
\end{itemize}

\textbf{Mechanism implication}

Sixflags does not fit a simple ``low task diversity'' story. A better
source-based description is:

\begin{itemize}
\tightlist
\item
  the workflow library is \textbf{narrow and park-centric}
\item
  the eval tasks are \textbf{mixed}, including both park/ticket flows
  and out-of-family information tasks
\item
  baseline CLICK is already relatively strong, so the narrow workflow
  family adds little headroom and can become harmful when it overfires
\end{itemize}

\subsubsection{Comparative summary
table}

{\def\LTcaptype{} % do not increment counter
\begin{longtable}[]{@{}>{\raggedright\arraybackslash}p{0.16\linewidth}>{\raggedright\arraybackslash}p{0.24\linewidth}>{\raggedright\arraybackslash}p{0.18\linewidth}>{\raggedright\arraybackslash}p{0.18\linewidth}>{\raggedright\arraybackslash}p{0.20\linewidth}@{}}
\toprule\noalign{}
Site & WF specificity & Observed task-family breadth & WF-task alignment
& Source-based reading \\
\midrule\noalign{}
\endhead
\bottomrule\noalign{}
\endlastfoot
\texttt{kayak} & Low & Medium & Strong & Abstract travel-search
primitives line up with a coherent search-oriented task family \\
\texttt{newegg} & Medium-Low & Medium & Strong & Commerce workflows are
still reusable; site grammar is stable enough that both LM and Rule can
compete \\
\texttt{budget} & High & High & Weak & Workflow library mixes several
function-specific families; eval tasks span unrelated site functions \\
\texttt{sixflags} & Medium & Medium & Partial to Weak & Workflow library
is narrow and park-centric, but eval tasks are more mixed than the
current main-text label suggests \\
\end{longtable}
}

\subsubsection{What this appendix
supports}

This appendix strongly supports:

\begin{itemize}
\tightlist
\item
  \texttt{kayak} and \texttt{newegg} as sites where reusable workflow
  families align with observed task families
\item
  \texttt{budget} as the clearest heterogeneity/mismatch case
\item
  \texttt{sixflags} as a site where \textbf{narrow workflow family +
  limited headroom + partial mismatch} is a better explanation than
  simply ``low task diversity''
\end{itemize}

This appendix does \textbf{not} fully support:

\begin{itemize}
\tightlist
\item
  a strong claim that \texttt{sixflags} has low task diversity
\item
  a strong claim that \texttt{WF\ specificity} alone explains
  success/failure
\end{itemize}



% C5
\clearpage
\subsection{Mind2Web Compositionality Reading}\label{app:C5}

\subsubsection{Overview}

The AWM paper uses workflow composition as part of its broader
narrative, especially in more open-ended settings. On Mind2Web, the
question is subtler:

\begin{itemize}
\tightlist
\item
  do the induced workflows show any sign of reusable subflows?
\item
  if yes, is that sign closer to true composition, or only to a flat
  library of reusable routines?
\end{itemize}

This note answers that question by reading the actual workflow texts and
contrasting stronger vs weaker cases.

This appendix is based on the audited workflow texts for
\texttt{kayak}, \texttt{united}, \texttt{newegg}, and \texttt{budget},
read together with the workflow-text evidence in
Appendix~\ref{app:C1}.

\subsubsection{Three Senses of
``Compositionality''}

To avoid over-claiming, this appendix separates three different notions:

\begin{enumerate}
\def\labelenumi{\arabic{enumi}.}
\item
  \textbf{Explicit hierarchical composition}

  \begin{itemize}
  \tightlist
  \item
    one workflow directly calls or nests another
  \item
    this is \textbf{not observed} in the current Mind2Web workflow texts
  \end{itemize}
\item
  \textbf{Flat subflow reuse}

  \begin{itemize}
  \tightlist
  \item
    a task can be completed by chaining several short reusable routines
  \item
    this is \textbf{clearly present on some sites}
  \end{itemize}
\item
  \textbf{Template bundling}

  \begin{itemize}
  \tightlist
  \item
    workflows are reusable only in a weak sense because they remain
    long, site- or task-specific bundles
  \item
    this is also present, especially on weaker sites
  \end{itemize}
\end{enumerate}

The current Mind2Web evidence supports (2) much more than (1).

\subsubsection{Stronger Evidence of Flat Subflow
Reuse}

\paragraph{\texorpdfstring{\texttt{kayak}}{kayak}}

The \texttt{kayak} workflow libraries provide the cleanest evidence of
subflow-style reuse.

In the \texttt{kayak} \texttt{LM} workflow library,
the workflow set is split into short units such as:

\begin{itemize}
\tightlist
\item
  \texttt{enter\_search\_location}
\item
  \texttt{select\_calendar\_date}
\item
  \texttt{apply\_amenity\_filter}
\item
  \texttt{sort\_search\_results}
\item
  \texttt{set\_flight\_trip\_type}
\item
  \texttt{adjust\_guests\_rooms}
\item
  \texttt{book\_selected\_option}
\end{itemize}

These are not end-to-end task scripts. They are local routines that can
plausibly be chained in different orders depending on the task:

\begin{itemize}
\tightlist
\item
  enter location
\item
  select date
\item
  adjust guest count
\item
  sort/filter
\item
  select a result
\end{itemize}

The corresponding \texttt{kayak} offline workflow library shows the same
pattern in a slightly broader form:

\begin{itemize}
\tightlist
\item
  location entry
\item
  travel date selection
\item
  flight type selection
\item
  filter selection
\item
  deal viewing
\end{itemize}

This is the strongest Mind2Web evidence that AWM is not merely
memorizing full trajectories. On \texttt{kayak}, the workflow library
really looks like a set of reusable subroutines.

\paragraph{\texorpdfstring{\texttt{united}}{united}}

\texttt{united} also shows substantial flat compositionality.

In the \texttt{united} \texttt{LM} workflow library,
the workflow set decomposes tasks into separable units:

\begin{itemize}
\tightlist
\item
  \texttt{enter\_flight\_locations}
\item
  \texttt{select\_travel\_dates}
\item
  \texttt{navigate\_travel\_info}
\item
  \texttt{search\_rental\_car}
\item
  \texttt{enter\_trip\_credentials}
\end{itemize}

Again, these are not one-shot whole-task solutions. They correspond to
pieces of larger airline-site tasks:

\begin{itemize}
\tightlist
\item
  flight search setup
\item
  date picking
\item
  account/trip retrieval
\item
  travel-information navigation
\item
  car-rental branching
\end{itemize}

So \texttt{united}, like \texttt{kayak}, supports a library-style
reading of composition: not nested function composition, but
recombinable site routines.

\subsubsection{Weaker or More Ambiguous
Cases}

\paragraph{\texorpdfstring{\texttt{newegg}}{newegg}}

\texttt{newegg} still shows reuse, but the compositionality signal is
weaker.

In the \texttt{newegg} \texttt{LM} workflow library,
the workflows include:

\begin{itemize}
\tightlist
\item
  \texttt{search\_product}
\item
  \texttt{apply\_price\_filter}
\item
  \texttt{apply\_selection\_filter}
\item
  \texttt{navigate\_category\_links}
\item
  \texttt{open\_site\_menu\_and\_select}
\end{itemize}

These are reusable, but they are also closer to \textbf{site utilities}
than to cleanly composable task fragments.

A task may use search + filter + category navigation, but the overall
structure is less obviously decomposed into stable subflows than on
\texttt{kayak} or \texttt{united}. The workflows look reusable, yet the
library feels shallower and more page-local.

\paragraph{\texorpdfstring{\texttt{budget}}{budget}}

\texttt{budget} is the clearest counterexample to a strong
compositionality story.

The \texttt{budget} offline workflow library mixes:

\begin{itemize}
\tightlist
\item
  car-rental location search
\item
  vehicle/extras selection
\item
  date/time entry
\item
  reservation lookup
\item
  career-site job search
\item
  deals browsing
\end{itemize}

Some of these routines are internally structured, but the library as a
whole is heterogeneous and weakly unified. It is better described as a
\textbf{bundle of extracted templates} than as a coherent compositional
vocabulary.

This matters because it explains why ``workflow reuse'' on Mind2Web is
highly site-dependent. A site can contain repeated sequences without
supporting a strong composition story.

\subsubsection{Online Workflow
Compression}

The online workflows weaken explicit compositionality even further.

For example, the \texttt{kayak} online workflow library contains
routines such as:

\begin{itemize}
\tightlist
\item
  \texttt{search\_for\_cars\_kayak}
\item
  \texttt{select\_location\_kayak}
\item
  \texttt{select\_dates\_for\_rental\_kayak}
\item
  \texttt{browse\_car\_type\_kayak}
\end{itemize}

These are useful, but more trajectory-shaped and narrower than the LM or
offline libraries. So even when they are reusable, they look more like
compressed routines tied to a recent successful path than like a broad
compositional basis.

This is consistent with Appendix~\ref{app:A7}:
online memory is closer to local test behavior, but less robust as a
general reusable library.

\subsubsection{Bottom-Line Reading}

The current Mind2Web evidence supports the following layered conclusion:

\begin{itemize}
\tightlist
\item
  \textbf{explicit hierarchical composition:} not supported
\item
  \textbf{flat subflow reuse:} clearly supported on \texttt{kayak} and
  \texttt{united}
\item
  \textbf{weaker, utility-style reuse:} present on \texttt{newegg}
\item
  \textbf{template bundling rather than true composition:} visible on
  \texttt{budget}
\end{itemize}

So the fairest critical-reproduction statement is:

On Mind2Web, AWM does show signs of subflow reuse, but mostly as a flat
library of short reusable routines rather than as explicit hierarchical
workflow composition. This makes the paper\textquotesingle s composition
narrative partially visible, but much weaker and flatter than the
stronger composition story associated with more open-ended environments.

\subsubsection{Interpretation limits}

\begin{itemize}
\tightlist
\item
  This appendix is a qualitative reading of the audited workflow
  libraries, not a formal graph-structure analysis.
\item
  It should support qualitative discussion of composition, not be turned
  into a hard quantitative claim.
\end{itemize}
